% GENERATED FILE. Edit canonical lessons, then run npm run generate:books.
% canonical-source-hash: fnv1a64:583a07d65526ec8e
% canonical-lessons: GU-C360-iman, GU-C360-lalac, GU-C360-vaphadar, GU-C360-dago, GU-C360-phasavvu

\chapter{Loyalty and Betrayal}
\label{ch:gu-loyalty-and-betrayal}
\hlchaptermodality{360}

\begin{chapteropening}
\textbf{By the end of this chapter:} \emph{I can talk about honesty, temptation and betrayal.}

Complete the last lesson of chapter 360: I can talk about honesty, temptation and betrayal.
\end{chapteropening}

\section[\texorpdfstring{\gu{ઈમાન} (īmān)}{ઈમાન (īmān)}]{\gu{ઈમાન} (īmān) — integrity}

\label{lesson:GU-C360-iman}



\begin{quote}
Before the new one: say the Gujarati for conceit, then the Gujarati for to bow.
\end{quote}

\subsection*{You'll want to know: \gu{ઈમાન}}
\textbf{\gu{ઈમાન}} — \emph{īmān} — \textquotedblleft{}integrity\textquotedblright{}.

\begin{quote}
\gu{સાચો} \gu{માણસ} \gu{પૈસા} \gu{માટે} \gu{પોતાનું} \gu{ઈમાન} \gu{વેચતો} \gu{નથી}.
\end{quote}

\begin{quote}
\emph{sāco māṇas paisā māṭe potānũ īmān vecto nathī.}
\end{quote}

\begin{quote}
An honest person does not sell his integrity for money.
\end{quote}

A neuter noun.

\begin{cousinweb}
\textbf{\gu{ઈમાન}} (\emph{īmān}) is usually traced to Arabic \emph{īmān}, \textquotedblleft{}faith\textquotedblright{}, through Persian; in Gujarati it is the faith one keeps with others.
\end{cousinweb}

\begin{grammarlens}[title={What you've built}]
One more word for giving reasons and opinions.
\end{grammarlens}

\subsection*{Your turn}
\begin{itemize}
\raggedright
  \item \emph{Say it:} \emph{īmān}
  \item \emph{Say it:} \emph{īmān}, once more
  \item \emph{Say it:} \emph{sāco māṇas paisā māṭe potānũ īmān vecto nathī.} — An honest person does not sell his integrity for money.
  \item \emph{Recall:} say the Gujarati for conceit, then the Gujarati for to bow, then say \emph{īmān} again
\end{itemize}

\begin{tcolorbox}[breakable,colback=teal!4,colframe=teal!35!black,title={Before you move on}]
What does \gu{ઈમાન} mean? (\textquotedblleft{}Integrity\textquotedblright{}.) Say it once more.
\end{tcolorbox}
\section[\texorpdfstring{\gu{લાલચ} (lālac)}{લાલચ (lālac)}]{\gu{લાલચ} (lālac) — temptation}

\label{lesson:GU-C360-lalac}



\begin{quote}
Before the new one: say the Gujarati for to bow, then the Gujarati for integrity.
\end{quote}

\subsection*{You'll want to know: \gu{લાલચ}}
\textbf{\gu{લાલચ}} — \emph{lālac} — \textquotedblleft{}temptation\textquotedblright{}.

\begin{quote}
\gu{તે} \gu{પૈસાની} \gu{લાલચમાં} \gu{ફસાઈ} \gu{ગયો}.
\end{quote}

\begin{quote}
\emph{te paisānī lālacmā̃ phasāī gayo.}
\end{quote}

\begin{quote}
He was lured by the money.
\end{quote}

A feminine noun.

\begin{cousinweb}
\textbf{\gu{લાલચ}} (\emph{lālac}) is usually traced to Sanskrit \emph{lālasa}, \textquotedblleft{}eagerly longing for\textquotedblright{} (CDIAL 11026), though CDIAL records no Gujarati form.
\end{cousinweb}

\begin{grammarlens}[title={What you've built}]
One more word for giving reasons and opinions.
\end{grammarlens}

\subsection*{Your turn}
\begin{itemize}
\raggedright
  \item \emph{Say it:} \emph{lālac}
  \item \emph{Say it:} \emph{lālac}, once more
  \item \emph{Say it:} \emph{te paisānī lālacmā̃ phasāī gayo.} — He was lured by the money.
  \item \emph{Recall:} say the Gujarati for to bow, then the Gujarati for integrity, then say \emph{lālac} again
\end{itemize}

\begin{tcolorbox}[breakable,colback=teal!4,colframe=teal!35!black,title={Before you move on}]
What does \gu{લાલચ} mean? (\textquotedblleft{}Temptation\textquotedblright{}.) Say it once more.
\end{tcolorbox}
\section[\texorpdfstring{\gu{વફાદાર} (vaphādār)}{વફાદાર (vaphādār)}]{\gu{વફાદાર} (vaphādār) — loyal}

\label{lesson:GU-C360-vaphadar}



\begin{quote}
Before the new one: say the Gujarati for integrity, then the Gujarati for temptation.
\end{quote}

\subsection*{You'll want to know: \gu{વફાદાર}}
\textbf{\gu{વફાદાર}} — \emph{vaphādār} — \textquotedblleft{}loyal\textquotedblright{}.

\begin{quote}
\gu{કૂતરો} \gu{પોતાના} \gu{માલિકને} \gu{વફાદાર} \gu{રહે} \gu{છે}.
\end{quote}

\begin{quote}
\emph{kūtro potānā mālikne vaphādār rahe che.}
\end{quote}

\begin{quote}
A dog stays loyal to its owner.
\end{quote}

\begin{cousinweb}
\textbf{\gu{વફાદાર}} (\emph{vaphādār}) is usually traced to Persian \emph{vafādār}, \textquotedblleft{}faithful\textquotedblright{}, built from Arabic \emph{wafā}, \textquotedblleft{}faithfulness\textquotedblright{}, and the Persian \emph{-dār}, \textquotedblleft{}holding\textquotedblright{}.
\end{cousinweb}

\begin{grammarlens}[title={What you've built}]
One more word for giving reasons and opinions.
\end{grammarlens}

\subsection*{Your turn}
\begin{itemize}
\raggedright
  \item \emph{Say it:} \emph{vaphādār}
  \item \emph{Say it:} \emph{vaphādār}, once more
  \item \emph{Say it:} \emph{kūtro potānā mālikne vaphādār rahe che.} — A dog stays loyal to its owner.
  \item \emph{Recall:} say the Gujarati for integrity, then the Gujarati for temptation, then say \emph{vaphādār} again
\end{itemize}

\begin{tcolorbox}[breakable,colback=teal!4,colframe=teal!35!black,title={Before you move on}]
What does \gu{વફાદાર} mean? (\textquotedblleft{}Loyal\textquotedblright{}.) Say it once more.
\end{tcolorbox}
\section[\texorpdfstring{\gu{દગો} (dago)}{દગો (dago)}]{\gu{દગો} (dago) — a betrayal}

\label{lesson:GU-C360-dago}



\begin{quote}
Before the new one: say the Gujarati for temptation, then the Gujarati for loyal.
\end{quote}

\subsection*{You'll want to know: \gu{દગો}}
\textbf{\gu{દગો}} — \emph{dago} — \textquotedblleft{}a betrayal\textquotedblright{}.

\begin{quote}
\gu{જે} \gu{મિત્ર} \gu{પર} \gu{મને} \gu{વિશ્વાસ} \gu{હતો} \gu{તેણે} \gu{મને} \gu{દગો} \gu{આપ્યો}.
\end{quote}

\begin{quote}
\emph{je mitra par mane viśvās hato teṇe mane dago āpyo.}
\end{quote}

\begin{quote}
The friend I trusted betrayed me.
\end{quote}

A masculine noun.

\begin{cousinweb}
\textbf{\gu{દગો}} (\emph{dago}) is usually traced to Persian \emph{daghā}, \textquotedblleft{}deceit, fraud\textquotedblright{}.
\end{cousinweb}

\begin{grammarlens}[title={What you've built}]
One more word for giving reasons and opinions.
\end{grammarlens}

\subsection*{Your turn}
\begin{itemize}
\raggedright
  \item \emph{Say it:} \emph{dago}
  \item \emph{Say it:} \emph{dago}, once more
  \item \emph{Say it:} \emph{je mitra par mane viśvās hato teṇe mane dago āpyo.} — The friend I trusted betrayed me.
  \item \emph{Recall:} say the Gujarati for temptation, then the Gujarati for loyal, then say \emph{dago} again
\end{itemize}

\begin{tcolorbox}[breakable,colback=teal!4,colframe=teal!35!black,title={Before you move on}]
What does \gu{દગો} mean? (\textquotedblleft{}A betrayal\textquotedblright{}.) Say it once more.
\end{tcolorbox}
\section[\texorpdfstring{\gu{ફસાવવું} (phasāvvũ)}{ફસાવવું (phasāvvũ)}]{\gu{ફસાવવું} (phasāvvũ) — to entrap}

\label{lesson:GU-C360-phasavvu}



\begin{quote}
Before the new one: say the Gujarati for loyal, then the Gujarati for a betrayal.
\end{quote}

\subsection*{You'll want to know: \gu{ફસાવવું}}
\textbf{\gu{ફસાવવું}} — \emph{phasāvvũ} — \textquotedblleft{}to entrap\textquotedblright{}.

\begin{quote}
\gu{તે} \gu{નકલી} \gu{માલ} \gu{વેચીને} \gu{ગ્રાહકોને} \gu{ફસાવે} \gu{છે}.
\end{quote}

\begin{quote}
\emph{te naklī māl vecīne grāhkone phasāve che.}
\end{quote}

\begin{quote}
He traps customers by selling them fake goods.
\end{quote}

Past, with \gu{તેણે} (he or she …): \gu{તેણે} \gu{એક} \gu{ગ્રાહકને} \gu{ફસાવ્યો}, \emph{teṇe ek grāhakne phasāvyo}.

\begin{cousinweb}
\textbf{\gu{ફસાવવું}} (\emph{phasāvvũ}) is the causative partner of \textbf{\gu{ફસાવું}} (\emph{phasāvũ}), \textquotedblleft{}to get stuck\textquotedblright{}; CDIAL lists the Gujarati verb, \textquotedblleft{}to ensnare\textquotedblright{}, under Sanskrit \emph{spāśáyati}, \textquotedblleft{}causes to be bound\textquotedblright{} (CDIAL 13814).
\end{cousinweb}

\begin{grammarlens}[title={What you've built}]
One more word for giving reasons and opinions.
\end{grammarlens}

\subsection*{Your turn}
\begin{itemize}
\raggedright
  \item \emph{Say it:} \emph{phasāvvũ}
  \item \emph{Say it:} \emph{phasāvvũ}, once more
  \item \emph{Say it:} \emph{te naklī māl vecīne grāhkone phasāve che.} — He traps customers by selling them fake goods.
  \item \emph{Recall:} say the Gujarati for loyal, then the Gujarati for a betrayal, then say \emph{phasāvvũ} again
\end{itemize}

\begin{tcolorbox}[breakable,colback=teal!4,colframe=teal!35!black,title={Before you move on}]
What does \gu{ફસાવવું} mean? (\textquotedblleft{}To entrap\textquotedblright{}.) Say it once more.
\end{tcolorbox}
